\section{Contenerización con Docker}

La contenerización es una técnica de virtualización ligera que permite ejecutar
procesos aislados (\textit{contenedores}) compartiendo el kernel del sistema operativo
anfitrión, a diferencia de las máquinas virtuales que virtualizan el hardware completo.
Cada contenedor encapsula el sistema de ficheros, las dependencias y la configuración
de una aplicación, garantizando que su comportamiento sea idéntico independientemente
del entorno donde se ejecute.

\textbf{Docker} es la plataforma de contenerización más extendida. Define un formato
de imagen en capas (\textit{Dockerfile}) que permite describir de forma declarativa
y reproducible el entorno de ejecución de una aplicación. Las imágenes se almacenan
en registros (Docker Hub, registros privados) y pueden desplegarse en cualquier sistema
con Docker instalado sin modificaciones.

En el contexto de herramientas multimedia, la contenerización resuelve el problema
habitual de la gestión de dependencias: Docker permite empaquetar una aplicación junto con todas sus dependencias en una
imagen reproducible e independiente del sistema operativo anfitrión. En el contexto
de Voctomix, la contenerización resuelve uno de los principales obstáculos para su
adopción: la instalación manual de GStreamer y sus plugins en distintas distribuciones
Linux puede resultar tediosa y propensa a errores de versión, especialmente dado que
Voctomix requiere tanto los plugins básicos (\texttt{gstreamer1.0-plugins-base, -good})
como los opcionales y no libres (\texttt{-bad, -ugly, libav}).

Docker Compose extiende este concepto al nivel de sistema, describiendo en un único
fichero YAML el conjunto de servicios, redes y volúmenes que componen la instalación
completa. En Voctomix 2.0 se definen diez servicios organizados en una red bridge
privada (\texttt{voctomix\_net}): el núcleo (\texttt{voctocore}), las cuatro fuentes de
cámara, el generador de contenido de pausa y fuera de emisión (\texttt{stream\_blanker}),
el broker de mensajería (\texttt{rabbitmq}), el servicio de telemetría (\texttt{notario})
y el gestor de audio (\texttt{audio\_manager}).

Las fuentes de cámara comparten el \textit{network namespace} del contenedor
\texttt{voctocore} mediante la directiva \texttt{network\_mode: service:voctocore},
lo que les permite comunicarse por \texttt{localhost} sin overhead de red virtual y
simplifica la configuración de los puertos de entrada.

% --- Figura sugerida ---
% \begin{figure}[H]
%   \centering
%   \includegraphics[width=0.9\textwidth]{figuras/docker_compose.png}
%   \caption{Diagrama de servicios Docker Compose de Voctomix 2.0.}
%   \label{fig:docker_compose}
% \end{figure}
